\inicializarSeccion{Comparación de Variables Categóricas}

Con base a lo anterior, realizaremos una comparación entre las variables categóricas para identificar posibles diferencias significativas entre los grupos. Utilizaremos pruebas de independencia para evaluar si existen asociaciones significativas entre las variables.

Para la primera hipótesis, se había planteado que el nivel promedio de oxígeno disuelto en nuestra región es mayor al umbral mínimo óptimo establecido por las normativas ambientales (5 ml/L). Podemos asociar esta hipótesis con la profundidad del agua debido a que existe una alta correlación entre la profundidad y el nivel de oxígeno disuelto. 

Planteamos la prueba de independencia para evaluar si existe una asociación significativa entre la profundidad del agua y el nivel de oxígeno disuelto y obtenemos el siguiente resultado:

\begin{verbatim}
	Pearson's Chi-squared test

data:  table(datosEq1$Depth_zone, 
datosEq1$O2_cat)
X-squared = 15820, df = 3, p-value < 2.2e-16
\end{verbatim}

Lo cual indica que existe una asociación significativa entre la profundidad del agua y el nivel de oxígeno disuelto.

Para nuestra segunda hipótesis, se plantea que existe una asociación entre la zona de profundidad y la clase de densidad del agua. Planteamos la prueba de independencia y obtenemos el siguiente resultado:

\begin{verbatim}
	Pearson's Chi-squared test

data:  table(datosEq1$Depth_zone,
 datosEq1$Water_density_class)
X-squared = 13402, df = 6, p-value < 2.2e-16
\end{verbatim}

Lo cual indica que existe una asociación significativa entre la profundidad del agua y la densidad del agua.